% Source-derived chapter 03: Standard conjecture of Hodge type
% Original source: standard-conjectures-abelian-fourfolds
\chapter{Standard conjecture of Hodge type}
\label{standard-conjectures-abelian-fourfolds-chapter-03}
\source{standard-conjectures-abelian-fourfolds}

\begin{remark}[Source section]
\label{standard-conjectures-abelian-fourfolds-chapter-03-source}
\source{standard-conjectures-abelian-fourfolds}
\source{standard-conjectures-abelian-fourfolds}
This chapter reproduces the corresponding section of the retrieved
LaTeX source, including its definitions, statements, examples, and
proofs. It has not yet been translated into Lean.
\notready
\end{remark}

% The source section title was: Standard conjecture of Hodge type
In this section we recall some classical conjectures (due to Grothendieck) and give some reformulations of those. We state our main result (Theorem \ref{thmscht}) and prove a consequence (Theorem \ref{thmclozel}).

Throughout the section, $k$ is a base field, $H^*$ is a fixed classical Weil cohomology and $R$ is the associated realization functor (see Conventions). We fix a variety $X$  over $k$ of dimension $d$ and  an ample divisor $L$ on $X$. We will write $\mathfrak{h}(X)$ for its homological motive (with respect to $H^*$). The section is written under the Assumption \ref{kun+lef} (which is satisfied in particular by abelian varieties).

\begin{conj}
\source{standard-conjectures-abelian-fourfolds}
\notready\label{conjkun} (Standard conjecture of K\"unneth type)
\source{standard-conjectures-abelian-fourfolds}


\noindent There exists a decomposition 
\[\mathfrak{h}(X)=\bigoplus_{n=0}^{2d} \mathfrak{h}^n(X)\] 
such that $R(\mathfrak{h}^n(X))=H^n(X)$ and $h^n(X)=h^{2d-n}(X)^{\vee}(-d).$
\end{conj}

\begin{conj}
\source{standard-conjectures-abelian-fourfolds}
\notready\label{conjlef} (Standard conjecture of Lefschetz type)
\source{standard-conjectures-abelian-fourfolds}

\noindent For all $n \leq d$, there is an isomorphism
\[\mathfrak{h}^n(X) \isocan \mathfrak{h}^{2d-n}(X)(d-n)\]
induced by the cup product with $L^{d-n}$. Moreover, the K\"unneth decomposition can be refined to a  decomposition 
\[\mathfrak{h}^n(X)= \mathfrak{h}^{n,\prim}(X)\oplus \mathfrak{h}^{n-2}(X)(-1)\] 
 which realizes to the primitive decomposition of $H^n(X)$.
\end{conj}

\begin{remark}
\source{standard-conjectures-abelian-fourfolds}
\notready
By the work of Lieberman and Kleiman the two conjectures above are known for abelian varieties \cite{GKL,Lieb}, see also Theorem \ref{classic}.
\end{remark}
\begin{assumption}
\source{standard-conjectures-abelian-fourfolds}
\notready\label{kun+lef}
\source{standard-conjectures-abelian-fourfolds}
From now on we will assume that $X$ verifies the   Conjectures \ref{conjkun} and \ref{conjlef}.
\end{assumption}
\begin{remark}
\source{standard-conjectures-abelian-fourfolds}
\notready
Conjectures \ref{conjkun} and \ref{conjlef} induce an isomorphism
\[\mathfrak{h}^n(X) \isocan \mathfrak{h}^{n}(X)^{\vee}(-n).\]
By adjunction this gives a pairing
\[q_n : \mathfrak{h}^n(X) \otimes \mathfrak{h}^{n}(X)  \longrightarrow \one(-n).\]
By construction  the Lefschetz decomposition 
\[\mathfrak{h}^{n }(X)=\mathfrak{h}^{n,\prim}(X) \oplus \mathfrak{h}^{n-2}(X)(-1)\] 
is orthogonal with respect to this pairing.
\end{remark}
\begin{definition}
\source{standard-conjectures-abelian-fourfolds}
\notready\label{defpairingmot} For $n\leq d$, we define the pairing 
\source{standard-conjectures-abelian-fourfolds}
\[\langle \cdot, \cdot \rangle_{n,\mot} : \mathfrak{h}^n(X) \otimes \mathfrak{h}^{n}(X)  \longrightarrow \one(-n)\]
recursively on $n$, by slightly modifying the pairing $q_n$ of the above remark. We impose that  the Lefschetz decomposition 
\[\mathfrak{h}^{n }(X)=\mathfrak{h}^{n,\prim}(X) \oplus \mathfrak{h}^{n-2}(X)(-1)\] 
is still orthogonal, we impose the equality $\langle \cdot, \cdot \rangle_{n,\mot} =\langle \cdot, \cdot \rangle_{n-2,\mot}(-2) $ on $\mathfrak{h}^{n-2}(X)(-1)$  and finally on $\mathfrak{h}^{n,\prim}(X)$ we define \[\langle \cdot, \cdot \rangle_{n,\mot}  = (-1)^{n(n+1)/2}q_n. \]

 \end{definition}
 \begin{remark}
\source{standard-conjectures-abelian-fourfolds}
\notready\label{rempola}
\source{standard-conjectures-abelian-fourfolds}
If $k$ is embedded in $\C$, the Betti realization of $\langle \cdot, \cdot \rangle_{n,\mot}$ is a polarization of the Hodge structure $H^n(X)$.
\end{remark}
\begin{definition}
\source{standard-conjectures-abelian-fourfolds}
\notready\label{defpairing} We define  the space $\mathcal{Z}^{n,\prim}_{\num}(X)_{\Q} $  and the  pairing  $\langle \cdot, \cdot \rangle_{n}$.
\source{standard-conjectures-abelian-fourfolds}
\begin{enumerate}
\item For $n \leq d/2$, we  define the pairing
 \[\langle \cdot, \cdot \rangle_{n} : \mathcal{Z}^{n}_{\hom}(X)_{\Q}  \times  \mathcal{Z}^{n}_{\hom}(X)_{\Q}  \longrightarrow \Q=\End(\one)\]
as follows. For $\alpha, \beta  \in \mathcal{Z}^{n}_{\hom}(X)_{\Q}  = \Hom (\one, \mathfrak{h}^{2n}(X)(n))$ consider 
\[(\alpha \otimes \beta)  \in   \Hom (\one, \mathfrak{h}^{2n}(X) \otimes \mathfrak{h}^{2n}(X)(2n))\] and define
\[\langle \alpha, \beta \rangle_{n} =      \langle \cdot, \cdot \rangle_{2n,\mot} (2n) \circ  (\alpha \otimes \beta).\]

\item For $n \leq d/2$, define  $\mathcal{Z}^{n,\prim}_{\num}(X)_{\Q} \subset \mathcal{Z}^{n}_{\num}(X)_{\Q}$ as
 \[\mathcal{Z}^{n,\prim}_{\num}(X)_{\Q}  = \Hom_{\NUM(k)_\Q}(\one, \mathfrak{h}^{2n,\prim}(X)(n)).\]  We will keep the same notation as in (1) for the induced pairing
 \[\langle \cdot, \cdot \rangle_{n} : \mathcal{Z}^{n,\prim}_{\num}(X)_{\Q}  \times  \mathcal{Z}^{n,\prim}_{\num}(X)_{\Q}  \longrightarrow \Q.\]
\end{enumerate}
 \end{definition}
 \begin{remark}
\source{standard-conjectures-abelian-fourfolds}
\notready
The  above definition is equivalent to the one given in the introduction, namely 
\[\mathcal{Z}^{n,\prim}_{\num}(X)_{\Q} = \{\alpha \in \mathcal{Z}^{n}_{\num}(X)_{\Q} , \hspace{0.2cm} \alpha \cdot L^ {{d-2n+1}}=0 \hspace{0.2cm}  \textrm{in}  \hspace{0.2cm}  \mathcal{Z}^{d-n+1}_{\num}(X)_{\Q}\}.\]
Moreover,  for $\alpha$ and $\beta$ in $\mathcal{Z}^{n,\prim}_{\num}(X)_{\Q}$,  we have the equality \[\langle \alpha, \beta \rangle_{n} =  (-1)^n \alpha \cdot \beta \cdot L^{d-2n}.\]
 Notice that in the introduction we
 did not make the Assumption \ref{kun+lef}. Instead Definition \ref{defpairingmot} cannot be formulated without the  Assumption \ref{kun+lef}.
\end{remark}
\begin{conj}
\source{standard-conjectures-abelian-fourfolds}
\notready\label{conjscht} (Standard conjecture of Hodge type)
\source{standard-conjectures-abelian-fourfolds}

\noindent For all $n \leq d/2$, the pairing  \[\langle \cdot, \cdot \rangle_{n} : \mathcal{Z}^{n,\prim}_{\num}(X)_{\Q}  \times  \mathcal{Z}^{n,\prim}_{\num}(X)_{\Q}  \longrightarrow \Q\]
of Definition \ref{defpairing} is positive definite.
\end{conj}


\begin{prop}
\source{standard-conjectures-abelian-fourfolds}
\notready\label{kerpairing}
\source{standard-conjectures-abelian-fourfolds}
The kernel of the pairing $\langle \cdot, \cdot \rangle_{n}$  on the vector space $\mathcal{Z}_{\hom}^{n}(X)_{\Q}$ is the space of algebraic classes  which are numerically trivial. The analogous statement holds true for $\mathcal{Z}_{\hom}^{n,\prim}(X)_{\Q}$.
\end{prop}
\begin{proof}

This is a direct consequence of a general fact. Let $T$ be a homological motive and suppose that it is endowed with an isomorphism
$f:T \isocan T^\vee$. Call $q: T \otimes T \rightarrow \one $ the pairing induced by adjunction. In analogy to Definition \ref{defpairing}, $q$ induces a pairing $\langle \cdot , \cdot \rangle$ on the space $\Hom(\one, T).$ First, notice that this pairing can also be described in the following way
\[\langle \alpha,\beta\rangle=\alpha^\vee \circ f \circ \beta.\]
This equality holds for formal reasons in any rigid category (alternatively, for homological motives, one can check it after realization).
Using this equality one has that $\beta$ is in the kernel of the pairing if and only if  $(\alpha^\vee \circ f) \circ \beta=0$ for all $\alpha$. As $f$ is an isomorphism, this is equivalent to the fact that $\gamma \circ \beta=0$ for all $\gamma\in\Hom(T,\one).$ This means precisely that $\beta$ is numerically trivial.

To conclude, notice that, by construction, one can apply this general fact to $T= \mathfrak{h}^{2n}(X)(n)$ or $T= \mathfrak{h}^{2n,\prim}(X)(n)$.
\end{proof}
\begin{corollary}
\source{standard-conjectures-abelian-fourfolds}
\notready\label{semipos} 
\source{standard-conjectures-abelian-fourfolds}
The pairing $\langle \cdot, \cdot \rangle_{n}$  is perfect on $\mathcal{Z}_{\num}^{n}(X)_{\Q}$. Moreover the pairing  is positive definite on  $\mathcal{Z}_{\num}^{n}(X)_{\Q}$ if and only if it is positive semidefinite on $\mathcal{Z}_{\hom}^{n}(X)_{\Q}$. The analogous statements hold true for $\mathcal{Z}_{\num}^{\prim,n}(X)_{\Q}$.
\end{corollary}
\begin{remark}
\source{standard-conjectures-abelian-fourfolds}
\notready\label{remHR}
\source{standard-conjectures-abelian-fourfolds}
If $k$ is embedded in $\C$ and if we work with Betti cohomology (cf. Remark \ref{rempola}), the Hodge--Riemann relations imply that $\langle \cdot, \cdot \rangle_{n}$  is positive definite on $(n,n)$-classes, hence on $\mathcal{Z}_{\hom}^{n}(X)_{\Q}$. In particular the standard conjecture of Hodge type holds true. 

Moreover, homological and numerical equivalence coincide (recall that we work under the Assumption  \ref{kun+lef}). The argument is due to Liebermann  \cite{Lieb} (see also \cite[Corollary 5.4.2.2]{Andmot}) and we recall it here.

If  $\alpha \in \mathcal{Z}_{\hom}^{n}(X)_{\Q}$ is a nonzero class  in codimension $n\leq d/2$ the inequality  $\langle \alpha, \alpha \rangle_{n}>0$ implies that $\alpha$ is not numerically trivial. On the other hand, the iterated intersection product with the hyperplane section $L$ induces an isomorphism $\mathcal{Z}_{\hom}^{n}(X)_{\Q}\cong \mathcal{Z}_{\hom}^{d-n}(X)_{\Q}$. Hence the intersection product $\mathcal{Z}_{\hom}^{n}(X)_{\Q}\times \mathcal{Z}_{\hom}^{d-n}(X)_{\Q}\rightarrow \Q$ is non-degenerate in one variable if and only if it is non-degenerate in the other, which implies that homological and numerical equivalence must coincide also in codimension $n\geq d/2$.

\end{remark}
Together with the case of characteristic zero (see the above remark), the following is the only known result on the standard conjecture of Hodge type \cite[5.3.2.3]{Andmot}.
\begin{theorem}
\source{standard-conjectures-abelian-fourfolds}
\notready
The pairing\label{segre} $\langle \cdot, \cdot \rangle_{1}$   on $\mathcal{Z}_{\num}^{1}(X)_{\Q}$ is positive definite. Equivalently, the pairing 
\source{standard-conjectures-abelian-fourfolds}
\[\mathcal{Z}^1_{\num}(X)_{\Q}\times \mathcal{Z}_{\num}^1(X)_{\Q} \longrightarrow \Q \hspace{1cm} D,D' \mapsto D\cdot D' \cdot L^{d-2}\]
is of signature  $ (s_+;s_-) = (1;  \rho_1-1)$ with
 $\rho_1=\dim \mathcal{Z}_{\num}^1(X)_{\Q}$.
\end{theorem}
\begin{prop}
\source{standard-conjectures-abelian-fourfolds}
\notready\label{anypolarization}
\source{standard-conjectures-abelian-fourfolds}
 %and $\tilde{\rho}_i=\sum_{j<i} \rho_j.$
The standard conjecture of Hodge type for a fourfold $X$  is equivalent to the statement that the intersection product
\[\mathcal{Z}^2_{\num}(X)_{\Q}\times \mathcal{Z}_{\num}^2(X)_{\Q} \longrightarrow \Q\]
is of signature  $ (s_+;s_-) = ( \rho_2 - \rho_1+1;  \rho_1-1)$ with
 $\rho_n=\dim \mathcal{Z}_{\num}^n(X)_{\Q}$.

In particular the conjecture does not depend on the ample divisor $L$ chosen.
\end{prop}

\begin{proof}
By Theorem \ref{segre} the standard conjecture of Hodge type holds true for divisors.  Hence we have to study codimension $2$ cycles.
Consider the decomposition \[\mathcal{Z}^{2}_{\num}(X)_{\Q} = \mathcal{Z}^{2,\prim}_{\num}(X)_{\Q} \oplus L \cdot \mathcal{Z}^{1}_{\num}(X)_{\Q}\] induced by the Lefschetz decomposition \[\mathfrak{h}^{4}(X)=\mathfrak{h}^{4,\prim}(X) \oplus \mathfrak{h}^{2}(X)(-1).\]  It is orthogonal with respect to the intersection product as the Lefschetz decomposition is orthogonal with respect to the cup product (already at homological level).

We claim that the intersection product on  $L \cdot \mathcal{Z}^{1}_{\num}(X)_{\Q}$ is of signature $(1;  \rho_1-1)$. Assuming the claim notice that  the intersection product is positive definite on $\mathcal{Z}^{2,\prim}_{\num}(X)_{\Q}$ if and only if the intersection product of the total space $\mathcal{Z}^2_{\num}(X)_{\Q}$ has the predicted signature.

For the claim, consider $\alpha=L \cdot D$ an element of $L \cdot \mathcal{Z}^{1}_{\num}(X)_{\Q}$. Then we have 
$\alpha\cdot \alpha = D \cdot D \cdot L^2$
hence the claim follows from Theorem \ref{segre}. 
\end{proof}
\begin{prop}
\source{standard-conjectures-abelian-fourfolds}
\notready \label{reducefinite}
\source{standard-conjectures-abelian-fourfolds}
Let $X_0$ be a specialization of $X$. Suppose that $X_0$ as well satisfies the Assumption \ref{kun+lef} (with respect to the specialization of $L$). Then the standard conjecture of Hodge type for $X_0$ implies the standard conjecture of Hodge type for $X$. 
\end{prop}
\begin{proof} Consider the canonical inclusion  $\mathcal{Z}_{\hom}^{n,\prim}(X)_{\Q}\subset \mathcal{Z}_{\hom}^{n,\prim}(X_0)_{\Q}$ induced by smooth proper base change in $\ell$-adic cohomology. Then, if the pairing $\langle \cdot, \cdot \rangle_{n}$  is positive semidefinite  on $\mathcal{Z}_{\hom}^{n,\prim}(X_0)_{\Q}$ it must be so for  $\mathcal{Z}_{\hom}^{n,\prim}(X)_{\Q} $. On the other hand, by Corollary \ref{semipos}, this semipositivity property is a reformulation of the standard conjecture of Hodge type.
\end{proof}
\begin{remark}
\source{standard-conjectures-abelian-fourfolds}
\notready
The above proposition reduces the standard conjecture of Hodge type to varieties defined over a finite field. Such reduction step is classical, see for example \cite[Remark 5.3.2.2(2)]{Andmot}.
\end{remark}

The following is our main result. It is shown at the end of Section 8.
\begin{theorem}
\source{standard-conjectures-abelian-fourfolds}
\notready\label{thmscht}
\source{standard-conjectures-abelian-fourfolds}
The standard conjecture of Hodge type holds for abelian fourfolds in positive characteristic.
\end{theorem}

\begin{corollary}
\source{standard-conjectures-abelian-fourfolds}
\notready\label{thmhomnum}
\source{standard-conjectures-abelian-fourfolds}
Let $A$ and $A_0$ be two   abelian fourfolds and suppose that  $A_0$ is a specialization of $A$. Let us fix a prime number $\ell$. If  $\ell$-adic homological equivalence on $A_0$ coincide with numerical equivalence on $A_0$ then the same holds true for $A$.
\end{corollary}
\begin{proof} First, notice that the question whether homological and numerical equivalence coincide only matters for $\mathcal{Z}_{\hom}^{2,\prim}(A)_{\Q}$. Indeed, homological and numerical equivalence coincide for divisors \cite{matsu} (see also \cite[Proposition 3.4.6.1]{Andmot}). This implies   that they coincide also on dimension one cycles,   as the standard conjecture of Lefschetz type holds true for abelian varieties. Finally, consider the decomposition \[\mathcal{Z}_{\hom}^{2 }(A)_{\Q} =\mathcal{Z}_{\hom}^{2,\prim}(A)_{\Q} \oplus L \cdot \mathcal{Z}_{\hom}^{1}(A)_{\Q},\]
and notice that on the complement of $\mathcal{Z}_{\hom}^{2,\prim}(A)_{\Q} $ the equivalences again coincide, as a consequence of the case of divisors.

Now, by smooth proper base change we have $\mathcal{Z}_{\hom}^{2,\prim}(A)_{\Q}\subset \mathcal{Z}_{\hom}^{2,\prim}(A_0)_{\Q}$. If    homological and numerical equivalence  coincide on $A_0$ then the pairing $\langle \cdot, \cdot \rangle_{2}$ on  $\mathcal{Z}_{\hom}^{2,\prim}(A_0)_{\Q}$ is positive definite by Theorem \ref{thmscht}. Hence it is also positive definite on $\mathcal{Z}_{\hom}^{2,\prim}(A)_{\Q}$.   By  Proposition \ref{kerpairing}, this means that there are no nonzero algebraic classes in $\mathcal{Z}_{\hom}^{2,\prim}(A)_{\Q}$ which are numerically trivial.
\end{proof}
\begin{theorem}
\source{standard-conjectures-abelian-fourfolds}
\notready\label{thmclozel}
\source{standard-conjectures-abelian-fourfolds}
Let $A$ be an   abelian fourfold. Then numerical equivalence on $A$ coincides with $\ell$-adic homological equivalence on $A$ for infinitely many prime numbers $\ell$. 
\end{theorem}
\begin{proof}
When $A$ is defined over a finite field, this result is due to Clozel \cite{Clozel}. (Clozel's result actually holds true without the dimensional restriction.) We can reduce to the finite field case by Corollary \ref{thmhomnum}.
\end{proof}
\begin{remark}
\source{standard-conjectures-abelian-fourfolds}
\notready
The fact that homological and numerical equivalence should always coincide is also one of the four standard conjectures. The two others, namely K\"unneth and Lefschetz, being known for abelian varieties, Theorems \ref{thmscht} and \ref{thmclozel} imply that in order to fully understand the standard conjectures for abelian fourfolds what is missing is $\ell$-independency of $\ell$-adic homological equivalence.
\end{remark}
